\chapter{机器学习研究设计答案}

\begin{enumerate}[leftmargin=*]
  \item purge 按标签信息区间删除与验证区间相交的训练样本；embargo 是额外的时间缓冲，用于所采用划分下仍存在的相邻信息重叠。是否需要缓冲及其长度取决于具体信息区间，二者都不自动保证观察独立。
  \item 内层验证选择模型、特征和超参数；外层测试把这一整套选择程序当作黑箱估计其未来误差。若把外层反复拿来 early stopping 或挑窗口，外层也变成选择数据，必须另设未见确认区间或降低证据等级。
  \item 训练标签最晚在 $\max I_{\rm train}+h=7+2=9$ 完成，早于验证起点 $10$；验证标签最晚在 $\max I_{\rm val}+h=11+2=13$ 完成，早于最终测试起点 $14$。本例两次使用的 2 都是标签跨度 $h$；额外缓冲应另行说明。
  \item Brier 分数 $\mathrm{BS}=n^{-1}\sum(p_i-y_i)^2$ 只评价概率误差；校准改善可能让概率更可信，但不改变排序、仓位、换手和成本。答案应先给分数变化，再说明需要单独的收益和执行实验。
  \item 第一折只用训练索引 $0\!:\!3$ 预测 $4\!:\!5$，第二折只用 $0\!:\!5$ 预测 $6\!:\!7$，最后按原时间顺序拼接预测。每次拟合只使用当时已经实现的标签，这样每个预测都来自没有拟合其自身标签的模型。
  \item 对每个 $k$ 取 top-$k$ 重要性集合 $A_k,B_k$，计算 $J_k=|A_k\cap B_k|/|A_k\cup B_k|$。同时报告不同随机种子/窗口的 $J_k$；若只在 $k=2$ 稳定、$k=20$ 崩溃，结论必须限定为短名单稳定。
  \item 测试集参与 early stopping 后，它已经影响模型选择，不能再叫外层样本外。修复是保留未见确认集，或把所有调参视为已消耗测试并重新采集后移数据；答案报告证据强度降级而不是只换标签。
  \item 输入分布改变时，先区分真实变化与测量方式变化；标签基准概率改变时，概率校准可能有帮助；条件预测关系改变时，则需要重新考察模型。成交或成本变化即使不影响预测，也会改变净收益。重新校准和重新估计都应在新的评价数据上比较，无法仅由一项漂移统计量决定最好的调整。
\end{enumerate}

\subsection*{基础衔接：独立计算}
多 $(0.5-0.3)^2=0.04$；两者期望损失分别为 $0.25,0.21$。
